\section{Q4：研究经验如何改善研究者？}
\label{sec:learning}

科研系统可以改进当前方案，却未必更善于开展下一项研究。Q4 关注这一区别：一次研究结束后留下什么信息，改变哪个组件，又能否帮助后续决策。我们区分任务内适应（within-task adaptation）、相关任务之间的跨运行复用，以及向留出任务族的迁移；它们是不同评价设置，不能互换。持久数据库能够支持前两者，却未必证明第三者；单个问题中的参数更新也可能完全局限于当前任务。科研品味还需要额外检验哪些问题值得投入，而不只是更高效地解决给定问题。

\subsection{外部记忆、可复用技能与局部反思}

AgentRxiv 让后续科研智能体读取既往研究报告；Agora 用 Git 记录共享假设、结果及核验关系~\cite{agentrxiv,agora}。两者把科研视为持续积累的集体过程，区别于普通对话历史的关键是跨会话持久化与复用。然而，持久化本身不能识别记忆的因果贡献：后续智能体也可能只是获得更好的候选解、更多累计计算或更充分的任务内证据。Agora 明确指出，单位算力的发现收益仍需匹配对照。

EvoScientist 区分构思记忆与实验记忆，从轨迹中提炼研究方向和实现策略；InternAgent-1.5 则组合程序化策略、轨迹内实验片段与跨会话语义知识~\cite{evoscientist,internagent-1-5}。这些架构启发我们区分记忆中的\emph{发生了什么}、\emph{为何选择该行动}和\emph{经验何时适用}。这是本文的设计解释，并非已经验证的通用记忆格式。AutoResearchClaw 强调从先前失败中提炼经验，其跨运行消融主要支持避免已知失败模式，而非提高科学质量上限~\cite{autoresearchclaw}。

可复用技能（reusable skills）不仅存储报告，还描述在可识别条件下应采取的行动。HypoForge 根据比较式批评和执行反馈修订假设生成、检验技能，在不更新基础模型的情况下评价留出任务~\cite{hypoforge}。这种策略适应有研究价值，但其检验通过率包含参考假设达到显著性的要求，不能单独确立假阳性控制或选题重要性。技能可能学会满足评价器，却仍不擅长产生有信息量的科学证据。

局部反馈也应保留独立定位。Reflexion 无需参数更新即可保存语言反思；Dolphin 将实验分析反馈给下一轮；OR-Agent 在启发式搜索中使用分层反思与回溯~\cite{reflexion,dolphin,or-agent}。它们为根据反馈调整行动提供基础，但不能把各自的实证范围改写为通用跨任务科研学习。尤其是 OR-Agent 的语言梯度、动量只是对反思机制的类比，并非对语言模型进行梯度更新。

\subsection{改变科研程序，还是训练模型？}

代码演化改变产生后续行动的过程。ADAS 根据既有设计与评价档案搜索智能体程序；Darwin G\"odel Machine 维护能够修改自身实现的编程智能体档案~\cite{adas,darwin-godel-machine}。二者为科研智能体设计提供基础，但证据来自给定基准，而非科学议程选择。Bilevel Autoresearch 进一步把内层科研循环本身作为优化对象：外层分析轨迹，生成改变搜索方式的机制~\cite{bilevel-autoresearch}。其报告实验限于一个小型 GPT 预训练场景，展示的是有界的外层干预内层，而非无限递归自我改进或通用跨领域迁移。

参数训练改变提出或评价决策的模型。CycleResearcher 利用自动审稿器开展研究生成的迭代偏好训练~\cite{cycleresearcher}。其核心评价包括模型生成但未执行的实验结果，附录另展示了与执行模块结合的案例。因此，稿件或审稿分数改善不能解释为实证发现能力改善。对训练数据而言，这一区别同样关键：流畅的事后解释未必记录了一场真正发生的实验。

ResearchTrails 从人类 Git 历史提取决策序列，研究示例、技能提炼、监督学习与强化学习~\cite{researchtrails}。其主要预测目标是记录中的下一步科研决策。提交历史只是人类探索的不完整证据，与历史决策相似也不能证明另一种策略会获得更多发现。SciSense 则合成基于引用文献的 sensemaking 轨迹，包括朝向已知贡献的路径重建，并将其蒸馏进规划模型~\cite{scisense}。这种重建能够提供结构化监督，却不是作者真实思考过程的观测；其下游执行实验与判断型评价也仍窄于通用科学价值。

MLE-Ideator 提供更直接的执行信号：构思模型通过强化学习学习策略指导，由另一智能体实现建议~\cite{mle-ideator}。留出的 MLE 任务检验指导能力能否迁移，但目标仍由外部给定。TTT-Discover 在单个问题的搜索中更新模型参数，Evolution Fine-Tuning 则先从进化轨迹训练，再评价留出的优化任务~\cite{ttt-discover,evolution-fine-tuning}。这一对照区分了任务局部适应与跨任务经验训练；EFT 的给定目标和主要采用的共同搜索框架，仍限制了对陌生科研工作流及 taste 的推断。

\begin{table}[tb]
\centering
\small
\caption{经验改变什么？机制与迁移是不同问题，同一工作可以采用多种机制。}
\label{tab:learning-mechanisms}
\begin{tabularx}{\linewidth}{@{}P{0.18\linewidth}YY@{}}
\toprule
机制 & 改变的对象与时间范围 & 代表性工作 \\
\midrule
局部反思 & 一次研究内的上下文与行动选择 & Reflexion；Dolphin；OR-Agent \\
记忆与技能 & 后续会话复用的外部证据或程序 & AgentRxiv；Agora；EvoScientist；InternAgent-1.5；AutoResearchClaw；HypoForge \\
程序演化 & 控制器代码或科研循环策略 & ADAS；Darwin G\"odel Machine；Bilevel Autoresearch \\
参数训练 & 任务前或任务中更新的模型权重 & CycleResearcher；ResearchTrails；SciSense；MLE-Ideator；TTT-Discover；EFT \\
\bottomrule
\end{tabularx}
\end{table}

\subsection{现有证据的限制与开放问题}

贯穿这些机制的核心难点是归因。更好的最终产物可能来自更多尝试、检索到的答案、更强的执行器，也可能来自更好的决策策略。因此，评价需说明保留了什么，尽可能固定执行接口，同时比较相同在线预算与包含经验采集的总成本；任务间是否重置记忆、权重也必须明确。如果训练和测试共享数据集、代码仓库或高度相似的任务模板，仅划分轨迹并不足够。

由此产生三个问题：什么经验能保留决策理由和适用条件，而不是复制成功答案？环境改变后，研究者能否修订曾经有效的策略？即使最终基准分数没有提升，学习能否改善信息性实验的选择和停止决策？回答这些问题需要真实执行与独立核验，不能只依赖同时提供训练奖励的评审器。下述两个 pilot 均为拟议设计，并非已报告结果；任务数和预算是有意控制规模的起步假设。

\subsection{拟议方案 A：识别科研记忆的因果价值}

\paragraph{假设与方法。}
同时保存成功和失败尝试的结构化决策记录，比润色后的报告摘要更可靠地改善后续实验选择。建立固定档案，记录状态、候选行动、当时理由、实际结果与尚未解决的不确定性。区分实现错误、证据不足和反驳假设的证据，避免统称为失败。将同一档案转换为文字摘要或结构化记录，保留来源链接；目标评价前冻结档案，防止不同条件相互污染。

\paragraph{数据、划分与对照。}
作为 pilot，拟从表格预测、图像分类和小型语言模型训练中选择十二个使用公开数据的小任务。每个任务族分配两个任务构建档案、一个开发、一个作为隐藏目标，底层数据集及数据划分在这些角色间互不重叠。先确认任务可运行，再预先登记任务身份。在相同目标与随机种子上随机分配六个条件：无记忆、完整文字摘要、仅成功记录、仅失败记录、完整决策记录、打乱任务对应关系的完整记录。摘要与完整记录使用相同底层证据，成功与失败对照用于隔离内容贡献；非空条件统一检索策略与记忆 token 上限。打乱的是任务与记忆的对应关系，而非记录顺序，同时保持数量和长度，以检验相关经验是否重要。

\paragraph{指标与资源核算。}
每个目标任务运行三个种子，仅按验证数据选择最终产物，记录其经独立核验的测试结果、重复失败率、冗余实验、无效结论，以及达到预先指定验证改进所需的时间。对抽样决策日志开展盲审，检查实验是否针对当时的不确定性。统一硬件，每次最多六个实验、45 分钟任务计算；如果每次始终占用一张设备，主要目标对照的上限为 40.5 device-hours。每个来源任务采集两次运行，每个开发任务及条件运行一次，分别增加九个和 13.5 个 device-hours 的上限。摘要、检索、模型 token 与失败执行另行核算，并报告实际设备利用率和 API 成本。这些设备时间上限不含远程模型计算，后者仍须计入成本。

\paragraph{失败风险。}
相似任务可能把记忆变成答案缓存；结果标签可能把低检验力误写成负发现；长记录也可能只是通过额外上下文产生帮助。打乱对照、来源检查和匹配证据只能部分缓解这些问题。目标任务仅有三个，区间应按任务聚类，并作为探索性结果解释。如果无收益能够区分防错帮助与科研决策改善，同样具有研究价值。

\subsection{拟议方案 B：跨留出任务族蒸馏科研策略}

\paragraph{假设与方法。}
从决策状态、备选行动、实际结果及适用条件训练，比仅模仿下一步行动更能迁移。固定实现智能体与评价器，仅改变策略模型。先用小型开放模型进行参数高效监督微调，待监督信号显示收益后再考虑强化学习。每个训练结果必须链接真实执行日志；未执行的合成理由可以作为注释保留，不能充当观测或行动成功的证据。

\paragraph{数据、划分与比较。}
拟构建十六个任务，表格预测、图像分类、小型语言模型训练和数值算法优化四个任务族各四个。按整个任务族留出：两个用于训练、一个开发、一个测试；每个任务内部再分离模型拟合数据、可见验证反馈与封闭最终评价。初始 pilot 只采用一个预先登记的任务族划分，通用迁移结论还需要进一步轮换。每个训练任务最多采集六条短轨迹，每条包含四次实验决策，合计最多 192 个实际执行的决策转移。比较基础策略模型、从同一档案检索的基础模型、仅行动微调、完整决策上下文微调；使用相同决策转移样本与优化步数预算，报告 token 数并在开发集增加匹配 token 的对照，再消融结果信息或适用条件。

\paragraph{指标、成本与解释。}
在留出任务族上冻结全部学习组件，四个目标任务各运行三个种子，每次最多四个实验、30 分钟；单设备执行时，四种方法的目标对照上限为 24 device-hours。监督微调的初始上限设为已声明硬件上的八个 device-hours。教师轨迹、筛选、开发、训练和推理分别核算，同时报告成本分摊到不同数量未来任务后的表现。主要指标是独立核验的任务改进及单位总预算收益；次要指标包括无效实验、过早停止和对评价器接口变化的敏感性。训练样本少、任务族惯例及数据污染都可能制造或掩盖迁移。即使结果为正，它支持的也只是给定目标下可复用的实验策略，而非自主选择重要科研问题。
